\documentclass[12pt]{article}
\usepackage[utf8]{inputenc}
\usepackage[T1]{fontenc}
\usepackage{hyperref}
\title{Kohoroxitari Language Resources for K-12 Curriculum Development}
\author{Compiled by Claude Sonnet 5.5 (claude-sonnet-5-5) for Andr\'as Kornai}
\date{}
\begin{document}
\maketitle

\section{Introduction}
Kohoroxitari (formerly ISO 639-3 kob) is not a separate language. Glottolog lists it as a retired entry, a village name, merged into Sanuma (xsu) in 2007; the village's speakers use the Yanoma (possibly Xitari) dialect of Sanuma. Resources should be sought under Sanuma.

\section{Primary Educational Resources}
None found under this name.

\section{Linguistic and Reference Materials}
\begin{itemize}
\item Glottolog entry for Kohoroxitari (retired): \url{https://glottolog.org/resource/languoid/id/koho1245}
\item ISO 639-3 record for kob: \url{https://publications.rifs-potsdam.de/cone/iso639-3/resource/kob}
\end{itemize}

\section{Religious and Community Resources}
None found.

\section{Digital and Audio Resources}
None found.

\section{Curriculum Development Context}
No curriculum can be built for Kohoroxitari as such; the relevant work is that for Sanuma, which was not researched here.

\end{document}